\section{Structured Abstract and Cumulative Proposal Map}
\label{app:abstract}
\begin{table}[H]
\caption{Eight-row structured abstract; actual results and future work remain distinct.}
\begin{tabularx}{\textwidth}{@{}p{0.28\textwidth}Y@{}}
\toprule
Proposal element & Current statement / change\\\midrule
Background & Collective decisions may benefit from information acquired privately before action; costly acquisition creates a public-good incentive question.\\
Motivation and application & In an AI-assisted group decision, agents or designers must decide who bears research cost when all benefit. Users of collective recommendations are stakeholders.\\
Three connected questions & Economics asks how cost affects free-riding; computation asks whether equilibria and cost changes are reproducible; behavior asks whether choices reflect strategic incentives or bounded beliefs. Together: when will agents acquire information, and will observed choices match the benchmark?\\
Method and model assumptions & Two agents simultaneously choose Research or Skip; one Research supplies value \(V\) to both and each researcher pays \(c\). Baseline \(V=4,c=2\), common knowledge, Nash equilibrium, NashPy plus independent best responses.\\
Results or expected results & \textbf{Actual:} two pure equilibria, symmetric \(p(R)=0.5\) mixed equilibrium, verified interior sweep \(0.75,0.50,0.25\), passing tests, executed notebook, and public Static Space. \textbf{Planned:} AI or human behavioral tests of choices and beliefs; no behavioral result exists.\\
Intellectual merits & The reproducible game clarifies a cost-sharing mechanism and specifies a behavioral gap that incentive theory alone cannot settle. Economics, computation, and behavior play distinct roles.\\
Practical impacts & \textbf{Existing:} a public interactive prototype for inspecting payoffs and equilibria. \textbf{Potential:} after validation, research-responsibility designs could inform organizational and public decisions. \textbf{Not evaluated:} improved learning or institutional efficiency. A pre/post best-response task could test the SDG 4 learning aim.\\
Feedback received and revision made & V1 focused on attention/interruption scheduling. The author reassessed its insufficiently explicit strategic interdependence and rebuilt v2 around costly information acquisition, aligning economics, computation, and behavior around one mechanism. One accessible review concerns v1; another assigned review is currently inaccessible. Neither is treated as direct v2 evidence.\\
\bottomrule
\end{tabularx}
\end{table}
